% figures/rvv-register-file.tex
% Diagram: RVV vector register file (v0–v31), each VLEN bits wide.
\tikzsetnextfilename{ch3-riscv-rvv-register-file}
\begin{tikzpicture}[
  reg/.style={draw, minimum width=6cm, minimum height=0.35cm,
              fill=blue!8, font=\ttfamily\scriptsize},
  label/.style={font=\ttfamily\scriptsize, anchor=east},
  annotation/.style={font=\sffamily\scriptsize, text=gray!70!black}
]

% Draw registers v0–v31 (show subset for clarity)
\foreach \i in {0,1,2,3} {
  \node[reg] (v\i) at (0, -\i*0.5) {};
  \node[label] at ([xshift=-3mm]v\i.west) {v\i};
}

% Ellipsis
\node[annotation] at (0, -2.25) {$\vdots$};

% v28–v31
\foreach [count=\idx from 0] \i in {28,29,30,31} {
  \pgfmathsetmacro{\ypos}{-2.75 - \idx*0.5}
  \node[reg] (v\i) at (0, \ypos) {};
  \node[label] at ([xshift=-3mm]v\i.west) {v\i};
}

% VLEN annotation
\draw[<->, thick, blue!60!black] ([yshift=2mm]v0.north west) -- ([yshift=2mm]v0.north east)
  node[midway, above, font=\sffamily\small] {\VLEN\ bits};

% Annotation: unknown at compile time
\node[annotation, anchor=west] at ([xshift=5mm]v0.east) {%
  \begin{tabular}{@{}l@{}}
    VLEN = hardware parameter\\
    (τυπικά 128–2048 bits)
  \end{tabular}};

% v0 special annotation
\draw[->, gray, thick] ([xshift=5mm]v0.east) ++(0,-1.2) -- ++(0,0.7)
  node[annotation, anchor=south west, xshift=1mm] {v0: mask register};

\end{tikzpicture}
